%% kapa: exam=1 round=29 year=2018 period=1 form=A subject=economics subject_ko=경제학원론 seq=25 number=65
%% source: Q-Net 감정평가사 기출문제 (cst00309) artlSeq=5202808 file=extracted/1008/5202808/2018년 제29회 감정평가사 1차시험 기출문제/2018년 제29회 감정평가사 1차시험 기출문제/1교시 A형.pdf
%% conversion: kordoc 4.24.51 -> tools/gen1.py (auto); review=proofread-sonnet; flags=none
\begin{question}{65}
소비이론에 관한 설명으로 옳은 것을 모두 고른 것은?
\begin{framedtext}
ㄱ. 케인즈 소비함수에 의하면 평균소비성향이 한계소비성향보다 크다. \par
ㄴ. 상대소득가설에 의하면 장기소비함수는 원점을 통과하는 직선으로 나타난다. \par
ㄷ. 항상소득가설에 의하면 항상소비는 평생 부(wealth)와 관계없이 결정된다. \par
ㄹ. 생애주기가설에 의하면 중년층 인구비중이 상승하면 국민저축률이 하락한다. \par
\end{framedtext}
\begin{choices}
  \item ㄱ, ㄴ
  \item ㄱ, ㄷ
  \item ㄴ, ㄷ
  \item ㄴ, ㄹ
  \item ㄷ, ㄹ
\end{choices}
\end{question}
\answer{1}
